% GENERATED, edit the French unit or the translation table
\mxnewpage\section*{Solutions EV4-B}\label{EV4-B-CORR}
\textbf{MCQ:} 1 A, 2 B, 3 C, 4 C, 5 D, 6 B, 7 C, 8 C, 9 D, 10 C, 11 A, 12 D. Each correct answer: 0.5 point; total 6.
\paragraph{Checks of the fluency questions.} Order of operations; sum of fractions; subtraction of exponents; multiplying coefficient 1.2; division by 0.8; product 1.25 times 0.8; solving an affine equation; vertical increment divided by horizontal; mean; angle divided by 360; complement to 1; substitution in brackets.
\subsection*{Exercise 1 -- An average price}
\begin{enumerate}
\item Values 0{,}4,6; probabilities 0.3; 0.3; 0.4. \textit{Marking: 1 point. Result and basic justification expected; half a point possible for a correct but unfinished approach.}
\item $E(X)=3.6=\dfrac{18}{5}$; $E(X^2)=19.2=\dfrac{96}{5}$; variance $\dfrac{156}{25}$. \textit{Marking: 2 points. Method or justification: 1; calculation and conclusion: 1.}
\item Mean $\dfrac{13}{5}$; standard deviation $\dfrac{2\sqrt {39}}{5}$, unchanged. The price -1 would be a refund: check whether this policy is really intended. \textit{Marking: 2 points. Method or justification: 1; calculation and conclusion: 1.}
\end{enumerate}
\subsection*{Exercise 2 -- Exponential maximum}
\begin{enumerate}
\item $f\prime=t(2-t)e^{-t}$. \textit{Marking: 1 point. Result and basic justification expected; half a point possible for a correct but unfinished approach.}
\item Increasing from 0 to 2 then decreasing; maximum $4e^{-2}$. \textit{Marking: 2 points. Method or justification: 1; calculation and conclusion: 1.}
\item $f(4)=16e^{-4}$; ratio to the maximum $4e^{-2}<1$ since $e^2>4$. \textit{Marking: 2 points. Method or justification: 1; calculation and conclusion: 1.}
\end{enumerate}
\subsection*{Exercise 3 -- Determining a position}
\begin{enumerate}
\item $(-x;-2)\cdot(4-x;-2)=x^2-4x+4=(x-2)^2$. \textit{Marking: 1 point. Result and basic justification expected; half a point possible for a correct but unfinished approach.}
\item Product zero only at x=2. \textit{Marking: 1 point. Result and basic justification expected; half a point possible for a correct but unfinished approach.}
\item Circle with centre (2;0), radius 2; the line y=2 is tangent to it at (2;2). \textit{Marking: 2 points. Method or justification: 1; calculation and conclusion: 1.}
\end{enumerate}
\paragraph{Calibration.} Target duration: 20--25 min for the fluency questions, 80--85 min for the exercises, 10--15 min for proofreading. These durations are editorial estimates to be measured in class. Accept any correct method; do not penalise twice a numerical error that is propagated consistently. The proposed marking scheme is not a decision of the team to harmonise marking.
